\documentclass[10pt,a4paper]{article}
\usepackage{amsmath,amssymb,amsthm,graphicx}
\usepackage{amsmath,amssymb,amsthm,hyperref,xfrac}
\providecommand{\eg}{e.g.\ }
\providecommand{\ie}{i.e.\ }
\usepackage[scr=rsfs,cal=euler]{mathalfa}
\multlinegap0pt
\hyphenation{accu-rate applied assump-tion assump-tions decom-pos-ing defined esti-mate esti-mates exten-ded hypothe-ses para-me-ter para-me-ters remain-der required}
\newenvironment{enumeratei}
{\bgroup\def\theenumi{\roman{enumi}}\def\theenumii{\arabic{enumii}}\begin{enumerate}}
{\end{enumerate}\egroup}
\newenvironment{enumeratea}
{\bgroup\def\theenumi{\alph{enumi}}\begin{enumerate}}
{\end{enumerate}\egroup}

\newcommand{\Psfrac}[2]{(\sfrac{#1}{#2})}
\newcommand{\psfrac}[2]{\sfrac{(#1)}{#2}}
\newcommand{\spfrac}[2]{\sfrac{#1}{(#2)}}
\newcommand{\pspfrac}[2]{\sfrac{(#1)}{(#2)}}
\newcommand{\shalf}{\sfrac{1}{2}}

\newcommand{\bnorm}[1]{\bigl\|#1\bigr\|}
\newcommand{\Bnorm}[1]{\Bigl\|#1\Bigr\|}
\newcommand{\babs}[1]{\bigl|#1\bigr|}
\newcommand{\Babs}[1]{\Bigl|#1\Bigr|}

\AtBeginDocument{%
\DeclareMathOperator{\sinn}{sin}\let\sin\sinn
\DeclareMathOperator{\coos}{cos}\let\cos\coos
}

%\usepackage{enumitem}
\usepackage{physics}
%\usepackage[foot]{amsaddr}
\usepackage{bm}
\usepackage{tikz}
\usepackage{mathtools}
\usepackage{scalerel,stackengine}
\stackMath
\usetikzlibrary{arrows.meta}
\usepackage{subcaption}
%\usepackage[style=numeric]{biblatex}
%\addbibresource{main.bib}
%\setlength{\bibitemsep}{0.3\baselineskip}

\newenvironment{mycases}
{\begin{dcases*}}
{\end{dcases*}}
\theoremstyle{plain}
\newtheorem{theorem}{Theorem}[section]
\newtheorem{lemma}[theorem]{Lemma}
\newtheorem{definition}[theorem]{Definition}
\newtheorem{proposition}[theorem]{Proposition}
\newtheorem{corollary}[theorem]{Corollary}
\newtheorem*{corollary*}{Corollary}
\theoremstyle{remark}
\newtheorem{remark}[theorem]{Remark}
\counterwithout{figure}{section}
\counterwithout{equation}{section}

\newcommand{\br}[1]{\left(#1\right)}
\newcommand{\bbr}[1]{\bigl(#1 \bigr)}
\newcommand{\Bbr}[1]{\Bigl(#1 \Bigr)}
\newcommand{\bbbr}[1]{\biggl(#1 \biggr)}
\newcommand{\tpd}{\frac{1}{(2\pi)^{d}}}
\newcommand\pig[1]{\scalerel*[2.35pt]{\big#1}{%
\ensurestackMath{\addstackgap[1.2pt]{\big#1}}}}
\newcommand\pigl[1]{\mathopen{\pig{#1}}}
\newcommand\pigr[1]{\mathclose{\pig{#1}}}
\newcommand{\ak}{\abs{k}}
\newcommand{\ap}{\abs{p}}
\newcommand{\aal}{\abs{\alpha}}
\newcommand{\sbbr}[1]{\bigl[#1\bigr]}
\newcommand{\sBbr}[1]{\Bigl[#1\Bigr]}
\newcommand{\sbbbr}[1]{\biggl[#1\biggr]}
\newcommand{\sbr}[1]{\left[#1\right]}
\newcommand{\half}{\frac{1}{2}}
\newcommand{\p}{\partial}
\newcommand{\jb}[1]{\langle#1\rangle}
\newcommand{\bset}[1]{\bigl\{#1\bigr\}}
\newcommand{\Bset}[1]{\Bigl\{#1\Bigr\}}
\newcommand{\set}[1]{\left\{#1\right\}}
\newcommand{\mytextfrac}[2]{{\textstyle\frac{#1}{#2}}}
\newcommand{\qaq}{\quad\text{and}\quad}
\newcommand{\para}{\parallel}
\newcommand{\ep}{\varepsilon}
\newcommand{\R}{\mathbb{R}}
\newcommand{\C}{\mathbb{C}}
\newcommand{\al}{\alpha}
\newcommand{\intRd}{\int_{\R^d}}
\newcommand{\intRdk}{\int_{\R^d_k}}
\newcommand{\intRdp}{\int_{\R^d_p}}
\newcommand{\lam}{\lambda}
\newcommand{\ceils}[1]{\lceil#1\rceil}
\newcommand{\Qin}{Q_{\mathrm{in}}}
\newcommand{\hin}{h_{\mathrm{in}}}
\newcommand\reallywidehat[1]{%
\savestack{\tmpbox}{\stretchto{%
\scaleto{%
\scalerel*[\widthof{\ensuremath{#1}}]{\kern.pt\mathchar"0362\kern.1pt}%
{\rule{0ex}{\textheight}}%WIDTH-LIMITED CIRCUMFLEX
}{\textheight}%
}{2.4ex}}%
\stackon[-6.9pt]{#1}{\tmpbox}%
}
\title{Fourier-space density decay implies decay of physical-space Sobolev norms}
\author{Marnie Smith}
\date{doi:10.5802/jep.353; CC BY4.0}
\begin{document}
\maketitle
\noindent\textit{Editorial note.} The counted statement is the author's unlabelled corollary environment at source0.tex lines 376-391; its proof environment opens at line 392, immediately after \textbackslash{}end\{corollary\}, so the pairing is by adjacency and the retained proof range 392-399 is the whole environment.\par\medskip
\noindent\textit{Editorial note.} The main theorem, with the labelled density-decay estimate T1 inside its statement (source lines 355-373), is retained as context because the counted proof applies that estimate at source line 398; the context is the complete theorem environment, which is also what supplies the label T1 referred to in the proof.\par\medskip
\noindent\textit{Editorial note.} Because the corollary carries no author label, its identity is recorded by the pinned source position of the two environments (statement lines 376-391, proof lines 392-399) and by the claim key, which encodes that position.\par\medskip
\noindent\textit{Editorial note.} This article carries no label-number dump and the retained counter values are left empty, so the printed numbering of the card is the wrapper's default; the original printed number of this unlabelled corollary is not recoverable from the pinned material and is therefore not asserted.\par\medskip
\noindent\textit{Editorial note.} The retained preamble is the publisher layout adapter, the native font and layout package and the native mathematical macro and theorem block of the pinned source.\par\medskip
\noindent\textit{Editorial note.} r002 note: the card is revived. r001 failed for two reasons that both sit on the card's side, not the author's: the material-preparation prose of this card, which the wrapper typesets, carried a literal Unicode sigma and nabla that the wrapper font does not provide, and the retained main-theorem context refers to the labelled Penrose condition display, which was not retained. r002 replaces those two characters in the card's own prose by equivalent TeX commands and retains that display as a further context. No retained author line is altered.\par\medskip
\noindent\textit{Editorial note.} r003 note: the retained material refers to the label NH at source line 227, which was not retained; r003 adds exactly that numbered environment as a further context and changes nothing else.\par\medskip
\noindent\textit{Editorial scope.} One original result of the pinned publisher article, counted once: physical-space decay of weighted Sobolev norms of the density (unlabelled corollary at source0.tex lines 376-391).\par\medskip
\noindent\textit{Allowed standard prerequisites (editorial).}\begin{itemize}
\item Fourier-space density decay. The Fourier transform of the density obeys a bound with a weight that grows in the frequency and in the product of frequency and time, with the exponent \$\textbackslash{}sigma\_1\$1 fixed by the hypotheses.. It is the estimate that the counted proof transfers to physical space.
\item Hausdorff-Young inequality. The Lp norm of a function is controlled by the Lp' norm of its Fourier transform, with 1/p + 1/p' = 1.. It is the first step of the proof.
\item Weighted derivative norms. The norm with the weight |\$\textbackslash{}nabla\$x, t\$\textbackslash{}nabla\$x|\textasciicircum{}n measures derivatives in space and in the time direction simultaneously.. It is the norm appearing in the conclusion and inside the integral.
\end{itemize}\medskip
\begin{theorem}[Nonlinear phase mixing]\label{maintheorem}
Fix dimension $d\geq3$ and $M:=\ceils{(d+1)/2}$. Choose constants $N_0>\sigma_1>\sigma_0$ satisfying \[\sigma_1\geq d+8,\quad N_0-\sigma_1>3d/2+7/2 \qaq \sigma_1-\sigma_0>d+1.\] Let $g\in H^{N_0+7/2+\delta_0}_{2+\ceils{d/2}}(\R^d;\R_+)$ for some constant $\delta_0\in(0,1]$, and let $w\in L^1(\R^d,\R)$ satisfy
\begin{align}
\abs{\widehat w(k)}\lesssim\frac{1}{\jb{k}^{M-1/2}}\label{wassumption}
\end{align}for all $k\in\R^d$. Assume $g$ and $w$ satisfy the uniform Penrose condition \eqref{Penrose} with constant ${c_{\mathrm{P}}}$. Then there exist constants $C>0$ and $\ep_0>0$, depending only on $N_0,\sigma_1,\sigma_0,w,g,{c_{\mathrm{P}}},d,\delta_0$, and independent of $\hbar\in(0,1]$, such that the following holds.

For any $\ep\in(0,\ep_0]$ and any self-adjoint initial data $\Qin\in\mathcal{H}^{N_0}_{M,M}$ with \[\norm{\Qin}_{\mathcal{H}^{N_0}_{M,M}}\leq\ep,\]
the nonlinear Hartree equation \eqref{NH} admits a unique global self-adjoint solution $Q(t)$, which satisfies:
\begin{itemize}
\item Density decay: for all $k\in\R^d$ and $t\geq 0$,\begin{align}
\label{T1}\abs{\widehat{\rho_Q}(t,k)}&\leq \frac{C\ep}{\jb{k, kt}^{\sigma_1}}.
\end{align}
\item Scattering: there exists an operator $Q_\infty\in\mathcal H^{\sigma_0}_M(\R^d;\C)$ such that, for all $t\geq 0$,
\begin{align}
\label{T2}
\norm{e^{-i\Psfrac{\hbar}{2}t\Delta}Q(t)e^{i\Psfrac{\hbar}{2}t\Delta}-Q_\infty}_{\mathcal H^{\sigma_0}_{M}}&\leq \frac{C\ep}{\jb{t}^{d/2}}.
\end{align}
\end{itemize}
\end{theorem}

\begin{align}\label{Penrose}
\inf_{\substack{\Re\lambda \geq 0 \\ k \in \R^d}} \abs{1 + \widetilde{\mathscr{L}}(\lambda, k)} \geq {c_{\mathrm{P}}}
\end{align}

\begin{align}\label{NH}
\begin{mycases}
i\hbar\p_t Q = \sBbr{-\frac{\hbar^2}{2}\Delta + w*\rho_{Q}, Q + \gamma_g}, \\
Q(0) = Q_{\mathrm{in}}.
\end{mycases}
\end{align}

\begin{corollary}[Physical-space decay]
Let $d\geq 3$, $p\in[2,\infty]$ and let $N_0>\sigma_1>0$ with
\begin{align*}
\sigma_1\geq d+8\qaq N_0-\sigma_1>3d/2+7/2.
\end{align*}
Then for any $n<\sigma_1-d(1-1/p)$, there exist constants $C>0$ and $\ep_0>0$, depending only on $d,p,n,N_0,\sigma_1$, and the parameters in Theorem~\ref{maintheorem}, and independent of $\hbar\in(0,1]$, such that the following holds.

For any $\ep\in(0,\ep_0]$ and any self-adjoint initial data $\Qin\in\mathcal{H}^{N_0}_{M,M}$ with
\begin{align*}
\norm{Q_{\mathrm{in}}}_{\mathcal H^{N_0}_{M,M}}\leq\ep,
\end{align*}
the solution $Q(t)$ to the nonlinear Hartree equation satisfies
\begin{align}\label{physicalspacedecay}
\norm{\jb{\nabla_x,t\nabla_x}^n\rho_Q(t,\cdot)}_{L^p(\R^d)}\leq \frac{C\ep}{\jb{ t}^{d(1-\sfrac{1}{p})}}.
\end{align}
\end{corollary}

\begin{proof}
By the Hausdorff--Young inequality,
\begin{align*}
\norm{\jb{\nabla_x,t\nabla_x}^n\rho_Q(t,\cdot)}_{L^p(\R^d)}
\lesssim_p \bbbr{\int_{\R^d} \jb{k,kt}^{p'n} \abs{\widehat{\rho_Q}(t,k)}^{p'} \,dk}^{1/p'}.
\end{align*}
Using the decay estimate \eqref{T1} and requiring convergence of the integral yields the result.
\end{proof}

\begin{thebibliography}{99}
\bibitem{r1} J. Bedrossian - “Nonlinear echoes and Landau damping with insufficient regularity”, Tunis. J. Math. 3 (2021) no. 1, p. 121-205
\bibitem{r2} J. Bedrossian, N. Masmoudi \& C. Mouhot - “Landau damping: paraproducts and Gevrey regularity”, Ann. PDE 2 (2016) no. 1, article ID 4, 71 pages
\bibitem{r3} J. Bedrossian, N. Masmoudi \& C. Mouhot - “Landau damping in finite regularity for unconfined systems with screened interactions”, Comm. Pure Appl. Math. 71 (2018) no. 3, p. 537-576
\bibitem{r4} J. Bedrossian, N. Masmoudi \& C. Mouhot - “Linearized wave-damping structure of Vlasov–Poisson in $\mathbb{R}^3$”, SIAM J. Math. Anal. 54 (2022) no. 4, p. 4379-4406
\bibitem{r5} A. Borie, S. Hadama \& J. Sabin - “Scattering for the positive density Hartree equation”, 2025
\bibitem{r6} A. Bove, G. Da Prato \& G. Fano - “An existence proof for the Hartree-Fock time-dependent problem with bounded two-body interaction”, Comm. Math. Phys. 37 (1974), p. 183-191
\bibitem{r7} A. Bove, G. Da Prato \& G. Fano - “On the Hartree-Fock time-dependent problem”, Comm. Math. Phys. 49 (1976) no. 1, p. 25-33
\bibitem{r8} T. Cazenave - Semilinear Schrödinger equations, Courant Lecture Notes in Math., vol. 10, American Mathematical Society, Providence, RI, 2003
\bibitem{r9} J. M. Chadam - “The time-dependent Hartree-Fock equations with Coulomb two-body interaction”, Comm. Math. Phys. 46 (1976) no. 2, p. 99-104
\bibitem{r10} T. Chen, Y. Hong \& N. Pavlović - “Global well-posedness of the NLS system for infinitely many fermions”, Arch. Rational Mech. Anal. 224 (2017) no. 1, p. 91-123
\bibitem{r11} T. Chen, Y. Hong \& N. Pavlović - “On the scattering problem for infinitely many fermions in dimensions $d\ge 3$ at positive temperature”, Ann. Inst. H. Poincaré C Anal. Non Linéaire 35 (2018) no. 2, p. 393-416
\bibitem{r12} J. J. Chong, L. Lafleche \& C. Saffirio - “On the $L^2$ rate of convergence in the limit from the Hartree to the Vlasov-Poisson equation”, J. Éc. polytech. Math. 10 (2023), p. 703-726
\bibitem{r13} C. Collot \& A.-S. de Suzzoni - “Stability of equilibria for a Hartree equation for random fields”, J. Math. Pures Appl. (9) 137 (2020), p. 70-100
\bibitem{r14} C. Collot \& A.-S. de Suzzoni - “Stability of steady states for Hartree and Schrödinger equations for infinitely many particles”, Ann. H. Lebesgue 5 (2022), p. 429-490
\bibitem{r15} E. Grenier, T. T. Nguyen \& I. Rodnianski - “Landau damping for analytic and Gevrey data”, Math. Res. Lett. 28 (2021) no. 6, p. 1679-1702
\bibitem{r16} S. Hadama - “Asymptotic stability of a wide class of stationary solutions for the Hartree and Schrödinger equations for infinitely many particles”, Ann. H. Lebesgue 8 (2025), p. 181-218
\bibitem{r17} A. D. Ionescu, B. Pausader, X. Wang \& K. Widmayer - “Nonlinear Landau damping for the Vlasov-Poisson system in $\mathbb{R}^3$: the Poisson equilibrium”, Ann. PDE 10 (2024) no. 1, article ID 2, 78 pages
\bibitem{r18} L. Landau - “On the vibrations of the electronic plasma”, Acad. Sci. USSR. J. Phys. 10 (1946), p. 25-34
\bibitem{r19} M. Lewin \& J. Sabin - “The Hartree equation for infinitely many particles, II: Dispersion and scattering in 2D”, Anal. PDE 7 (2014) no. 6, p. 1339-1363
\bibitem{r20} M. Lewin \& J. Sabin - “The Hartree equation for infinitely many particles I. Well-posedness theory”, Comm. Math. Phys. 334 (2015) no. 1, p. 117-170
\bibitem{r21} M. Lewin \& J. Sabin - “The Hartree and Vlasov equations at positive density”, Comm. Partial Differential Equations 45 (2020) no. 12, p. 1702-1754
\bibitem{r22} C. Mouhot \& C. Villani - “On Landau damping”, Acta Math. 207 (2011) no. 1, p. 29-201
\bibitem{r23} T. T. Nguyen \& C. You - “Modified scattering for long-range Hartree equations of infinite rank near vacuum”, SIAM J. Math. Anal. 57 (2025) no. 6, p. 6486-6497
\bibitem{r24} T. T. Nguyen \& C. You - “Plasmons for the Hartree equations with Coulomb interaction”, Probab. Math. Phys. 6 (2025) no. 3, p. 913-960
\bibitem{r25} E. M. Stein \& G. Weiss - “Interpolation of operators with change of measures”, Trans. Amer. Math. Soc. 87 (1958), p. 159-172
\bibitem{r26} C. You - “Phase mixing estimates for the nonlinear Hartree equation of infinite rank”, 2024
\end{thebibliography}
\end{document}
